%=========================================================
%   DOCUMENT PREAMBLE
%=========================================================

\documentclass[12pt,oneside]{article}

%---------------------------------------------------------
%   ENCODING, LANGUAGE, FONTS (pdflatex compatible)
%---------------------------------------------------------
\usepackage[T1]{fontenc}
\usepackage[utf8]{inputenc}
\usepackage[english]{babel}

%---------------------------------------------------------
%   PAGE LAYOUT
%---------------------------------------------------------
\usepackage{geometry}
\geometry{paper=a4paper, margin=1in}

\usepackage{setspace}
\onehalfspacing

%---------------------------------------------------------
%   HYPERLINKS
%---------------------------------------------------------
\usepackage{hyperref}
\hypersetup{
    colorlinks = true,
    linkcolor = blue,
    urlcolor  = blue,
    citecolor = blue,
    pdfauthor = {Philippe Thomas Savard},
    pdftitle  = {Philosophy Behind the Geometry of the Prime Number Spectrum}
}

%---------------------------------------------------------
%   USEFUL PACKAGES
%---------------------------------------------------------
\usepackage{amsmath, amssymb, amsthm}
\usepackage{graphicx}
\usepackage{enumitem}
\usepackage{listings}
\usepackage{csquotes}

%---------------------------------------------------------
%   SECTION STYLE
%---------------------------------------------------------
\setcounter{secnumdepth}{3}
\setcounter{tocdepth}{3}

%=========================================================
%   BEGIN DOCUMENT
%=========================================================

\begin{document}

%---------------------------------------------------------
%   TITLE
%---------------------------------------------------------
\begin{center}
{\Huge \textbf{Philosophy Behind the Geometry of the Prime Number Spectrum}}\\[1em]
{\Large by Philippe Thomas Savard}\\[2em]
\end{center}

%---------------------------------------------------------
%   TABLE OF CONTENTS
%---------------------------------------------------------
\tableofcontents
\newpage

\section{Autobiography: Academic Path and Early Experiences}

As I have mentioned several times throughout the different documents we produced together, I do not have a particularly complex or extensive academic background. When I was of age to pursue studies, I attempted to complete a college-level program in civil engineering at Cégep Beauce–Appalaches. It is a field where mathematics are omnipresent throughout the entire curriculum. I completed two years out of the three-year program, that is, four sessions out of six.

During that period—and this was true throughout my entire academic journey, from primary to secondary school and college—I consistently received poor grades in mathematics, except on two occasions. The first was during my final year of high school. In Secondary~5, I was placed in the most advanced mathematics class, known as Mathematics~536.

After the final ministry exam, my mathematics teacher came to see me during another class to tell me how proud he was of me, because I was the only student in the entire school who had passed the exam. The situation was unusual: without anyone realizing it, the teacher himself had not received the correct information from the ministry regarding the material that was supposed to be taught for the exam. As a result, the entire class failed. He was unsure whether a retake would be offered, which was unfortunate for the other students. He was genuinely happy that I had succeeded, according to what he told me—a moment I still remember clearly.

Surprised by the situation, I was deeply marked by what he said about my performance. Even though the Mathematics~536 curriculum had been taught to us, none of it appeared on the exam. I passed the exam without having been taught the material that was actually evaluated.

The truth is that I was never particularly strong in school—often inattentive, and frequently uninterested in the teaching being offered, even though I recognize its importance. But throughout my life, especially in mathematics, I always took what we learned in school and used it to do mathematics on my own, at home, for myself and my own questions. I often did mathematics homework—but my own homework, with the questions I wanted to answer.

And this, in my view, is why I succeeded on the ministry exam at the end of Secondary~5. Lacking the proper information to answer the questions, I relied on my personal mathematical intuition. I like to joke that I know how to “build all the numbers in the world,” and it is this ability that I applied during the exam.

\subsection{The College Experience}

In college, the same pattern repeated itself: I struggled in class, and it is no surprise that I eventually dropped out. But during one of the sessions, in one of the program’s courses, there was an exam in civil engineering that was intentionally designed by the teaching staff so that all students would fail. The purpose was to calibrate how strict they could be with the cohort when evaluating their understanding of the program material.

Once again, I passed this exam that everyone was expected to fail. Of course, the retake was successful for the rest of the class, since a retake was offered. Normally, exam questions in college follow the course outline, which students can consult throughout the session. But this time, the “rigged” exam was not based on the course outline at all—it was a general assessment of the field of civil engineering.

It was a kind of situational test, asking students to act as though they were already civil engineering technicians—not to answer questions about material taught in class, but rather: “How would you proceed if you were already a technician?” For me, the exam felt simple, because doing things and figuring out how to do them is something that comes naturally to me. I always have something to say about how things should be done—worldwide, even, I like to joke, though I exaggerate only slightly.

\subsection{Reflections on Others and the Life Drive}

Whenever I speak with people in my life—though it is often difficult to discuss this topic with them—I always wish to talk about a discussion within the discussion, something that fascinates me deeply. This meta-discussion concerns how people understand and resolve the personal difficulties they face in complex situations.

What I always want to ask them is: how did they do it? How did they manage to free themselves from those difficult situations? I often find that people are extremely harsh toward themselves. When they agree to talk about it, it is usually only after the situation has resolved itself somewhat, often through an unexpected event that saved them at the last moment. They are amazed by their luck and give themselves very little credit, sometimes even blaming themselves for not having been more careful.

To me, when people manage to escape such situations, it is because they are fundamentally good and capable individuals, with a structured sense of order in their lives. The order they maintain is proportional to the danger they must sometimes face. They are autonomous and resilient.

I know that what allows them to escape these situations is the life drive—the subtle force, the finesse. I define the life drive as: “the fantasy of the object that surpasses life through its reasons for being.” This is my own definition, inspired by another concept that intertwines with it: sadism, defined as “the fantasy of the object of the death drive.”

To understand my definition, one must see it as an expression of the spirit of finesse, which stands in contrast to the geometric spirit. The geometric spirit demands proofs and detailed demonstrations. The people around me who share how they escaped their difficult situations often describe their actions as hesitant or uncertain, even though they succeeded.

\section{Reflections on the Geometric Spirit, the Life Drive, and the Analogist}

When people recount these situations in which they were unexpectedly saved from difficulty, they often adopt the geometric spirit. I understand why: they do not want to encourage the very oversight in their personal management that led them into such precarious circumstances, from which they were rescued almost by chance. Because they are autonomous and capable individuals, skilled in organizing their lives, they refuse to believe that they could have committed an error that exposed them to such troubling possibilities.

Yet I am convinced of this: within every person’s existence lies everything necessary for them to find a way out of the complex situations they face.

The discourse of the fascist—who seeks to sever friendship, even though friendship is precisely the recognition that we do not have to do everything alone—spreads pamphlets of death. By this I mean that when we listen to such people speak about someone who is proud, or who believes they have succeeded, or who is openly facing their own success, these fascists always respond with the same refrains: “It’s because of your mother,” “It’s only because you have money,” “It’s because you have friends,” “It’s due to someone else’s knowledge,” “Your brother‑in‑law helped you,” “Your tools,” “Luck”… They distribute pamphlets of death, never granting the slightest merit to anyone.

The life drive—this finesse—when confronted with such discourse, can lead a fragile mind to believe that the only way to be right about one’s existence is to end it. Blinded by the idea that they are condemned to succeed without ever being the author of their success, and unable to recognize that the life drive always allows us to rise above circumstances using what already exists within our lives, they may come to believe that even their friends and family would agree with such a conclusion. They imagine that others, too, have reached the “truth” that they should disappear. Tragically, many answer this question with fatal acts far too often.

Mathematics have fascinated me for as long as I can remember. As mentioned earlier, I have always worked on mathematics personally, at every age and as far back as my memory goes. This is how I became a kind of grammarian of mathematics. An analogist is a grammarian who sees grammar as a form of mathematics.

Having checked, no university in Quebec teaches this discipline, so they cannot claim that I am not an analogist—one must first teach a subject before claiming expertise in it. Aristotle and Plato were analogists, and they founded the Academy of Knowledge, one of the earliest known universities in the world. They themselves had never attended such an institution.

An analogist is a grammarian who sees grammar as mathematics. Most grammatical facts are the work of analogists, since all numbers are written in letters. For example, 3 is written “three” and 2 is written “two,” and their numerical values are respectively 2 and 3 in terms of letter count.

However, the anomalist has long been considered more “true” than the analogist, because the anomalist sees grammar as causality, much like physics interprets the universe.

Synthetic reasoning is something that truly fascinates me, and I adopt it constantly. If every cause precedes its effect in time, then the cause must already know its effect. This is the approach I use daily, and I prefer it to causality, because anything reasoned in advance is far more human. I hunt at the grocery store, I eat from the refrigerator, and I sleep in a den surrounded by drywall: a natural environment for any wild Québécois.

An analogist is someone who intervenes when a form of know‑how becomes misleading. They verify and work to eliminate algorithmic biases that may affect society.

It is through what I call \textit{isossophy} that the analogist projects forward into the future, evaluating various subjects by comparing current values with those that are socially accepted. From this projection toward the future, the analogist observes whether these values remain in their proper place relative to other projected or already existing values—those that society wishes to preserve—while also maintaining a view toward our lived past.

It is a kind of geometric cohomology of impossible figures, where an object blocks our field of vision, yet through mathematical methods it becomes possible to deduce the object behind the one obstructing our view. But in the case of isossophy, the field of vision is truly blocked, and the only way to deduce what lies on the other side is through the hidden face of the very object that blocks our sight.

The word \textit{isossophy} is composed of two parts: “iso,” meaning equal measure, and “sophy,” which in ancient times referred to a specialist of knowledge. Today, however, a sophist is someone who argues in a deceptive or fallacious manner. The isossophist eliminates misleading and fallacious ways of functioning. Such methods or algorithms are recorded and placed in the annals and memory of society, and everything related to them thereafter is stripped of what these deceptive methods had fraudulently taken from the collective. This allows society to recover what was taken from it.

\section{Definition of Idioschizophrenia}

In earlier times, certain individuals whose behavior was questionable were noticed and recorded. In other words, their actions were documented, and these behaviors remain recognizable even today. This behavior was labeled without foundation and without any real concern for the common or collective good.

At a time when the general population had a very low level of literacy, some individuals took advantage of others’ ignorance to act without restraint and in dishonest ways. Coins already bore markings indicating their respective values—5 cents, 10 cents, 25 cents, and dollars. These dishonest individuals used deceptive strategies to pass off a 5‑cent coin as a 25‑cent coin or as another value, exploiting the population’s limited reading ability. Many citizens were ruined by such deceit. Their manipulative discourse was recorded and analyzed, and these analyses remain in use even today.

In the present day, as some individuals enter early adulthood, they are identified as engaging in discussions similar to those of the past. They are labeled as having a mental disorder—diagnosed, according to experts, with schizophrenia. They are handicapped not by law, but by society, which does not tolerate the quality of their discourse or their ideals. Their condition is acquired, meaning it is imposed upon them. Often placed under guardianship, they cannot sign contracts or payment agreements once their diagnosis becomes official and visible. They lose, quite literally, even their family ties and find themselves without an identity.

A notable element of their discourse, common among many of them at the onset of the disorder, is the belief that they are God. To the author, it seems understandable that someone who has lost their identity—who no longer remembers who they are—might come to believe they are the most well‑known figure in the world: the Good Lord.

These young adults, in their early twenties, reach this stage almost entirely thanks to caring parents. Once they reach adulthood, they sometimes clumsily claim, for example, that their father was “crazy” for having provided them with education, shelter, and resources that allowed them to reach the age of twenty. Yet they owe this far more to the fact that their father implemented structure and regulation in their lives—one of the main reasons for their success.

Often, while their parents continue to manage the obligations of their lives as they did when they were minors, these young adults declare, again clumsily, that their father was “crazy” for having raised them, even as they themselves struggle to meet their own obligations and life projects, often due to excessive social behavior. From society’s perspective, it is far more likely that it was the father’s structure—not his supposed madness—that allowed these young people to reach adulthood. Socially, this kind of discourse can have a significant impact on their future.

A well‑known definition of schizophrenia is: “Early disturbance of the mother‑child relationship. — Pierre Larousse Dictionary. Confused delirium.” This discourse of early derealization and depersonalization, typical at the beginning of adulthood, is recorded as being comparable to the dishonest individuals of earlier times, when low literacy levels made abuse easier. They are perceived as people seeking to exploit, in a morbid way, a form of social wealth. Rarely will a person diagnosed with schizophrenia, if they retain the right to manage their finances, achieve the financial stability and autonomy necessary for a balanced life.

“Doctus cum libro”: Said of those who are incapable of thinking for themselves, who display borrowed knowledge and draw their ideas from the works of others. (Definition from the Pierre Larousse Dictionary, pink pages: Doctus cum libro — “Learned with the book”.)

\subsection{Schizophrenia Diagnosed After Age 35}

A particular form of schizophrenia is the one that manifests after the age of 35. It is generally not diagnosed before a person has developed a clear understanding of their identity. To detect this late‑onset form, it is often necessary for a specialist to observe the individual clinically. Their discourse is highly destabilizing: speaking with someone diagnosed at this age can give the impression that it is the interlocutor—not the patient—who is delusional.

\subsubsection*{Fictional Account}

Two friends, a man and a woman, are walking near a cemetery. In the distance, they see a group of people gathered around a coffin, which they suspect belongs to one of their relatives. The group eventually leaves the burial site. At one moment, a young woman from the group stumbles and falls. The man walking with his friend bursts into laughter. His friend tells him:

“You know, you shouldn’t laugh. That woman who just fell is burying her father today, whom she loved very much. Every Sunday, she went horseback riding with him. Once, she fell off her horse, and that is what killed her father: everyone laughed at her for falling, and that is what destroyed him.”

This antinomy, where self‑reference is extremely present—“Once, she fell off her horse, and that is what killed her father”—illustrates, though fictitiously, the kind of self‑referential reasoning that a person with schizophrenia diagnosed after age 35 may exhibit.

When we face someone suffering from undiagnosed schizophrenia past the age of 35, we may feel as though we are the ones afflicted, if we rely on their machinations. One must understand what antinomy is: a form of reasoning that constantly seeks, in a morbid way, to select the element that cancels functioning. Often dishonest, antinomy relentlessly attempts to choose the position in a communication between two or more collaborators that invalidates or blocks the event, circumstance, project, emotion, feeling, phenomenon, etc.

\subsection{The Third Person Who “Wants”}

The self‑criticism of an individual suffering from idioschizophrenia reveals a trait of thought that overlays layers of reality. Each of us has our own reality, within which we organize ourselves to meet our needs. The individual suffering from the condition seeks to be a person in their imagination and thought, but not in tangible reality. Moreover, beyond wanting to be someone only in their imagination, the individual with idioschizophrenia seeks to feel this imagined person through a third party—that is, through someone other than who they are in tangible reality.

This characteristic of their self‑perception, in relation to those around them and to who they are relative to others, appears in their self‑criticism regarding the consequences and scope of their actions.

\section{Phenomenology of Idioschizophrenia}

The individual suffering from the condition justifies their behavior by referring to a third party who is not themselves, claiming that this third person would have had the will to act in one way or another, even before this third person had acted toward anyone or anything.

This psychophysical mechanism, in which their behavior is attributed to a third party unrelated to their reactions, serves to blame the consequence of their act on someone else, and additionally to accuse that third party of being responsible for their conduct due to this supposed will.

They thus construct an economy based on misconduct, where repayment must occur: the debt incurred by their actions must be repaid by the third party who supposedly “wanted” the individual with idioschizophrenia to act in such a way, even though this third party is often entirely unrelated to the victim or the person affected by their behavior.

Self‑reference, in their personal critique of this psychophysical action, reveals a chirality of reality from which, through this asymmetry, all toxic, deviant, delinquent, immoral, or improper behaviors emerge. This phenomenology of the individual suffering from idioschizophrenia regulates their cybernetics of action, justified by this third party who serves as scapegoat and mystery, allowing them in a second step to reproach the supposed will.

Cybernetics is the science of the mastered analogy between organism and machine. It is the science that studies communication. Machine as machination: they attempt to reduce to zero what terrifies them most in life—the good conduct of human beings.

Indeed, the dystopian side of idioschizophrenia is easily observed in their terror of young boys, toward whom they organize radical measures from infancy, referring to a dubious experience of illiteracy to justify the implementation of radical acts intended to prevent these young boys from one day having a wife—and that she might have breasts.

This unjustified fear, which in the present moment is already an impossibility at the age of these young individuals, is based on a questionable life experience meant to prevent something that is not terrifying in itself, by the very evidence of what it is: the obvious fact that an adult woman who is a partner will naturally have breasts.

This reality of the individual suffering from idioschizophrenia is the reason for my decision to classify it as a form of certain paranoid schizophrenia.

\subsection*{Idio: Analogy and Etymology}

“Idio” by analogy with idiom and idiosyncratic. An idiom is an expression, a word, a turn of phrase that is geographically localized to a given region. Outside the region where the idiom is used, it is not heard. Language allows idioms; idioms do not teach language how to speak.

Idiosyncratic refers to doing one thing and succeeding at something unexpected. It is the success of someone who is aware. A person who is not aware of what they are focusing on cannot achieve this kind of success, because it would go unnoticed.

Idioschizophrenia is localized to the status of the academic and is observed in conscious individuals who, by choice, transform success into the rank of the Doctor’s error. The syndrome of “The Specialist Physician.” Idioschizophrenia.

\section{Preamble: Knowledge}

Knowledge is more a memory than a novelty. It presents itself as a substance that is felt, like a chiasm: when we touch it, it touches us in return. Reminiscence and realism combine within it, similar to an acceleration. Three laws, modeled after Isaac Newton’s laws of acceleration, can be formulated to define what knowledge is.

\subsection{First Law: Consciousness}

For knowledge to truly exist, there must absolutely be consciousness. Without consciousness, no substance can be felt or recognized. The first law is therefore unavoidable and mandatory.

\subsection{Second Law: The Inverse of Knowledge — Understanding}

As mentioned in the first law, consciousness is an absolute value. But in the second law, the law of the inverse of knowledge, it is essential that consciousness be considered null.

At this precise position in time—the present moment—it becomes a reference frame: a kind of analogue relation to knowledge, a meta‑knowledge that manifests as an evident truth. Not knowing, and recognizing that we share a common sense with this absence of knowledge, constitutes a fixed point, a primitive of knowledge.

This fixed point guarantees the logic of knowledge and allows change, like a moment of inertia in which knowledge begins to move.

Thus, when the inertia of knowledge differs from zero—that is, when one recognizes not knowing—then the first law is indeed considered null.

Through the inverse of knowledge, our capacity to speak—understanding—emerges from this fixed point. The reference frame becomes animated and moves from the present toward the memory of past moments and circumstances. All the rules related to this forming knowledge manifest there: the capacity to speak must exceed our ignorance, otherwise nothing can truly be known.

\subsection{Third Law: Similar Figures}

As mentioned earlier, the rules related to knowledge are those that take shape. Knowledge is more a memory than a novelty. From past knowledge inscribed in our memory, we compare the rules of this forming knowledge with what is already known.

Through these comparisons, by adjusting our vocabulary to the event taking shape, and by relying on observable rules from the past within the established reference frame, we can access knowledge.

\subsection{2.1 Disproportioning What Is Known}

The expression “stubborn” is often used to describe a stubborn child. This child, who is in the process of learning to speak, cannot yet qualify a need, an object, a desire, or a necessity. Because they do not know, they call upon someone close to them—often a parent—and, in an attempt to qualify the element they cannot identify, they make a kind of absurd proposal to determine how to qualify what they struggle to name.

Persistent in their request, the child repeats the same absurd demand several times, and the parent usually refuses. You have seen the scene: eventually, the parent, exasperated, angrily grants the request. Immediately, the child, instead of enjoying the long‑requested opportunity, begins to cry and says they do not want it.

In Quebec, this behavior is also called “tête de cochon,” a stubborn streak typical of childhood, something we have all experienced and can still observe in the youngest, which is understandable since they do not yet fully know how to speak properly and cannot reasonably be expected to qualify everything at that age.

Idioschizophrenia is this early disturbance between the human and the learned human. It can be observed in individuals suffering from idioschizophrenia, just as in earlier times when people deceived others due to low literacy levels. Even today, some individuals still behave abusively.

Indeed, the “specialist physician syndrome” is an incapacity to know—one of the characteristic symptoms of this mental disorder. Only reproach allows them to manipulate an argument resembling that of someone who might know something.

Idioschizophrenia is an adult who seeks, through stubbornness, to disproportion the entirety of what is known so that it is no longer known, because the stunned interlocutor mistakenly believes they are the one responsible. The words of the afflicted person imply that the interlocutor is responsible for the misunderstanding.

Like solipsism, the goal of this paradoxical impression is to make the interlocutor doubt through the environment. For many philosophers, it is impossible to doubt through the environment.

For example, if there is a smell of gasoline, we do not doubt by saying: “If there were no gasoline fumes, I could light my lighter to check whether gasoline is in the air,” because lighting it would be madness.

The solipsism they invoke is an element of their delusion and dystopian logic. They invoke a past experience—often dubious—linked to their history, which in the present moment is already an impossibility, and this fear with questionable foundations justifies, in their view, the implementation of radical actions for a future they believe morbidly dangerous if such radical means are not applied immediately.

This frenzy, in which individuals suffering from idioschizophrenia weave the staging of a dangerous \emph{dasein} following a pseudo‑security discourse, reveals their condition when they proclaim themselves the police of what exists and what does not exist.

The authority bias is extremely present. According to them, their interlocutor’s statements “do not exist.” Let us call this passage “It‑does‑not‑exist‑to‑say.”

Then they justify all their behavior through this rule they impose in an unparalleled frenzy, a staging in which they evaluate everything negatively, simulating reality in order to disqualify and handicap in reality. Let us call this second passage “It‑is‑all‑because‑of‑to‑say.”

They are a secret; they carry a violent present, a pernicious wish they guard secretly in order to break something in an indelible way.

\section{Preamble: Knowledge}

Knowledge is more a memory than a novelty. It presents itself as a substance that is felt, like a chiasm: when we touch it, it touches us in return. Reminiscence and realism combine within it, similar to an acceleration. Three laws, modeled after Isaac Newton’s laws of acceleration, can be formulated to define what knowledge is.

\subsection{First Law: Consciousness}

For knowledge to truly exist, there must absolutely be consciousness. Without consciousness, no substance can be felt or recognized. The first law is therefore unavoidable and mandatory.

\subsection{Second Law: The Inverse of Knowledge — Understanding}

As mentioned in the first law, consciousness is an absolute value. But in the second law, the law of the inverse of knowledge, it is essential that consciousness be considered null.

At this precise position in time—the present moment—it becomes a reference frame: a kind of analogue relation to knowledge, a meta‑knowledge that manifests as an evident truth. Not knowing, and recognizing that we share a common sense with this absence of knowledge, constitutes a fixed point, a primitive of knowledge.

This fixed point guarantees the logic of knowledge and allows change, like a moment of inertia in which knowledge begins to move.

Thus, when the inertia of knowledge differs from zero—that is, when one recognizes not knowing—then the first law is indeed considered null.

Through the inverse of knowledge, our capacity to speak—understanding—emerges from this fixed point. The reference frame becomes animated and moves from the present toward the memory of past moments and circumstances. All the rules related to this forming knowledge manifest there: the capacity to speak must exceed our ignorance, otherwise nothing can truly be known.

\subsection{Third Law: Similar Figures}

As mentioned earlier, the rules related to knowledge are those that take shape. Knowledge is more a memory than a novelty. From past knowledge inscribed in our memory, we compare the rules of this forming knowledge with what is already known.

Through these comparisons, by adjusting our vocabulary to the event taking shape, and by relying on observable rules from the past within the established reference frame, we can access knowledge.

\subsection{Disproportioning What Is Known}

The expression “stubborn” is often used to describe a stubborn child. This child, who is in the process of learning to speak, cannot yet qualify a need, an object, a desire, or a necessity. Because they do not know, they call upon someone close to them—often a parent—and, in an attempt to qualify the element they cannot identify, they make a kind of absurd proposal to determine how to qualify what they struggle to name.

Persistent in their request, the child repeats the same absurd demand several times, and the parent usually refuses. You have seen the scene: eventually, the parent, exasperated, angrily grants the request. Immediately, the child, instead of enjoying the long‑requested opportunity, begins to cry and says they do not want it.

In Quebec, this behavior is also called “tête de cochon,” a stubborn streak typical of childhood, something we have all experienced and can still observe in the youngest, which is understandable since they do not yet fully know how to speak properly and cannot reasonably be expected to qualify everything at that age.

Idioschizophrenia is this early disturbance between the human and the learned human. It can be observed in individuals suffering from idioschizophrenia, just as in earlier times when people deceived others due to low literacy levels. Even today, some individuals still behave abusively.

Indeed, the “specialist physician syndrome” is an incapacity to know—one of the characteristic symptoms of this mental disorder. Only reproach allows them to manipulate an argument resembling that of someone who might know something.

Idioschizophrenia is an adult who seeks, through stubbornness, to disproportion the entirety of what is known so that it is no longer known, because the stunned interlocutor mistakenly believes they are the one responsible. The words of the afflicted person imply that the interlocutor is responsible for the misunderstanding.

Like solipsism, the goal of this paradoxical impression is to make the interlocutor doubt through the environment. For many philosophers, it is impossible to doubt through the environment.

For example, if there is a smell of gasoline, we do not doubt by saying: “If there were no gasoline fumes, I could light my lighter to check whether gasoline is in the air,” because lighting it would be madness.

The solipsism they invoke is an element of their delusion and dystopian logic. They invoke a past experience—often dubious—linked to their history, which in the present moment is already an impossibility, and this fear with questionable foundations justifies, in their view, the implementation of radical actions for a future they believe morbidly dangerous if such radical means are not applied immediately.

This frenzy, in which individuals suffering from idioschizophrenia weave the staging of a dangerous \emph{dasein} following a pseudo‑security discourse, reveals their condition when they proclaim themselves the police of what exists and what does not exist.

The authority bias is extremely present. According to them, their interlocutor’s statements “do not exist.” Let us call this passage “It‑does‑not‑exist‑to‑say.”

Then they justify all their behavior through this rule they impose in an unparalleled frenzy, a staging in which they evaluate everything negatively, simulating reality in order to disqualify and handicap in reality. Let us call this second passage “It‑is‑all‑because‑of‑to‑say.”

They are a secret; they carry a violent present, a pernicious wish they guard secretly in order to break something in an indelible way.

\section{Idioschizophrenia: Rupture Between Reality and Imagination}

The individual suffering from idioschizophrenia claims to possess a free spirit. This avowed antinomism manifests through the rejection of all responsibility in their rebellion, as though they had no shared sense of meaning with others, no empathy—those who do not suffer from this condition.

The refusal of all consciousness of propriety, agreement, peace, tranquility, law, social morality, family morality, legal morality, etc., reveals in them a troubling expansiveness of morality. Their discourse activates only when others make a decision: they then unfailingly choose to act in the opposite manner. They immediately wish to undertake radical measures to prevent the choices of others—who, in truth, never included them in their decision.

These individuals suffering from idioschizophrenia impose themselves and intrude, allowing themselves to criticize harshly and openly, even to the point of humiliating and insulting the choices and decisions of others.

For Philippe Thomas Savard, it is certain that the individual suffering from idioschizophrenia cannot access knowledge. Only through reproach do they attempt to approach knowledge, but without ever hoping to share even the slightest common meaning. In essence, they do not understand that they do not understand.

It is nearly impossible for these individuals to admit that they do not know something. And in their antinomic dialect, it is not uncommon for them to blurt out, without thinking, when asked to understand the value of their actions if those actions were directed toward them. Without consciousness, they often declare:

\begin{itemize}
    \item “I do not know what knowledge is.”
    \item “I do not know what understanding is.”
\end{itemize}

They have only one desire: to amuse themselves by believing that people they do not know at all must somehow know what they want—an attitude they share with the stubborn child.

To avoid having to criticize themselves, they refuse to acknowledge that people see with their own eyes, because they themselves do not know what seeing is.

These behaviors reveal a profound lack, a morbid obstinacy. For the author, this lack is the primary reason for their idioschizophrenia: the fact that they act and claim that it does not matter, and that the fact that someone does nothing to them would somehow have an impact on them.

Strange auditory phenomena multiply. It is not uncommon for a person suffering from the condition to claim that someone “knows” they said something, even though they never addressed that person, and to allege that this person should somehow know what they meant.

Their relationship with the facts of our tangible reality is primarily psychophysical, a mark of their intellectual impairment, and their self‑criticism—when they become aware of their actions—is mainly used to reproach others for events.

Their only connection to knowledge is always through a third person, and it is useful to them only for reproach. They want—or have—an opinion on a subject only once others have had that opinion. In most cases, what they want above all is a safe‑conduct to at least be able to reproach. For the author, this is their only connection to the possibility of knowing.

\subsection{Psychophysical Action}

In physics, every cause precedes its effect in time. It is certain that a black billiard ball struck by a white ball moves not because the white ball hits it, but because the exact moment the white ball collides with the black ball is the moment the black ball begins to move.

The individual suffering from idioschizophrenia excludes all a priori synthetic reasoning, claiming that “already” is a word that should be removed from the French dictionary—a change that, in their view, would make their behavior normal, since “the world” would no longer notice their improper actions because it would no longer remember.

Everything is reduced to banality, until we can easily observe their phenomenological perception (hallucination) of psychophysical action, which they judge and which reveals their loss of reference and contact with our tangible reality.

When someone does nothing to them, it affects them, and they complain—furiously—that this person, who has not even spoken to them, is doing something to them that is, according to them, atrocious and unbearable.

Conversely, when they commit a harmful act—even up to homicide—they claim that it “does not matter” and insist that they are the victims of the person in the present moment, even if that person is dead. They often claim that the deceased person is insulting them.

For the individual suffering from idioschizophrenia—the “specialist physician syndrome”—their actions are the consequence of a non‑present third party who supposedly “wants,” which justifies their often deadly misconduct toward someone other than this absent third party.

\subsection{Rupture Between Reality and Fiction}

The rupture between reality and fiction—imagination—is also profound in the individual suffering from idioschizophrenia. They assert vigorously and firmly that imagination and reality are identical, that there is no difference between the two.

They claim to have access to a technology that allows them to hear everything, to hear thoughts, and to speak within the thoughts of other individuals.

When the individual suffering from this condition elaborates on the subject, they claim that those with whom they communicate through this technology (hallucinatory) — even though they speak to them concretely — are in fact only hallucinations. They claim that these people, with whom they communicate without permission, merely believe they are communicating with them in their thoughts.

But for the individual suffering from idioschizophrenia, this imaginary communication is imperative: it must be true, it must be reality.

They even go so far as to confront the people they claim to have contacted through thought. And in person, they do not admit to having communicated with them originally: they diagnose them as suffering from deep hallucinations.

The level of harassment these individuals inflict on those they believe they have contacted through their so‑called technology is extremely serious.

They may go as far as planning criminal acts: homicide, rape, theft.

Most of them are dangerous offenders, often occupying positions at the university level—where access to personal and confidential information is particularly sensitive.

Their preferred places to act according to their “free spirit” are hospitals, since the most sensitive data are highly accessible there, and obviously very critical.

\section{Idioschizophrenia: Deep Depersonalization and Remarkable Rupture of Self and Other}

Idioschizophrenia manifests as a profound depersonalization and a marked rupture between self and other, a planned chirality that opposes and confuses the two. It is accompanied by a pronounced dream of “enslavement”: the observation of good conduct becomes a means of disinheriting the other from their best judgment, even as the individual rejects all forms of morality, which has deprived them of judgment.

The free spirit of antinomism studies good conduct—particularly that of men—to formulate a pernicious wish.

\subsection{Intelligence and the Distinction Between Real and Imaginary}

A person’s intelligence is a sequence of interpretations of symbols and signs.

\begin{itemize}
    \item Symbols affect the senses.
    \item Signs affect judgment.
\end{itemize}

It is a reading of symbols and signs that unfolds through synchronic and diachronic cycles. Intelligence is the tool that allows us to distinguish and differentiate the material world of reality, in which we live, from the immaterial world of imagination.

Without intelligence, only reality exists. Everything that belongs to knowledge and understanding—our ability to reason a priori about concepts, questions, situations, circumstances, etc.—does not exist.

For these individuals, the imaginary and all immaterial things are “the same as” material reality. Every fiction is assimilated to reality, and they accuse everything fictional of being real, without self‑criticism or remorse. A convenient way to perpetually harass and manipulate reality as they please.

\subsection{Depersonalization and Deviant Behaviors}

The individual suffering from idioschizophrenia seeks to be the person they are—but only in their thoughts. This person they attempt to become in their mind, they seek to feel through a second person who is not themselves.

From this depersonalization arise toxic, deviant, delinquent, or socially improper behaviors.

\subsection{The Fantasy of Enslavement}

The term “esclavagisme” is not accurate. The proper form is “esclavagiste,” which designates the person who engages in the trafficking of slaves. “Esclavage” is what the slave is subjected to.

There is no logical chain that leads someone to become a slave; this is why “esclavagisme” is not an expression in itself. “Esclavagisme” is more a fantasy than a reality.

The individual suffering from idioschizophrenia nourishes this fantasy intensely. They spend their time observing good conduct, particularly that of men. They watch for a term, a word, an expression, a turn of phrase that would allow them, through these elements, to subject these observed men to the slavery of a life of misery, misfortune, poverty, sadness, illness, death, etc.—forever.

This is the pernicious wish of the individual suffering from this condition: a secret handicap they wish to see imposed upon humanity in general.

This individual is misanthropic, meaning they harbor a general hatred of human beings. To illustrate this, the author imagines it as wanting to place one’s head inside one’s own mouth—a mystical image of hatred. This hatred is associated with the figure of death, sometimes represented with a large mouth.

Finally, the author describes idioschizophrenia as a condition resulting from an inability to know without reproach. The afflicted individual seeks to disproportion what is known so that this knowledge ceases to be recognized. This intellectual handicap can be heard when speaking with someone who suffers from it. They proclaim themselves the “police” of what exists and what does not exist, convinced they possess a superior sense of reality.

Quickly, they divert the conversation by asserting that what their interlocutor says “does not exist”—what the author calls the first response: “It‑does‑not‑exist‑to‑say.”

Then, they attribute everything the other person says to this initial assertion—the second response: “It‑is‑all‑because‑of‑to‑say.”

Such a discussion produces a morbid staging, a frenzy, in which the interlocutor feels discomfort. The afflicted person works to maraud the other’s opinion, to restrict the factors guiding their decisions and actions. They observe this process meticulously to provoke failures in the quality of the interlocutor’s judgment. The author calls this “frenzy.”

Thus, these individuals orchestrate a real‑time simulation to provoke the judgment error of others, depriving them of their best discernment in order to bring them down to their own handicap. In certain academic environments, such practices of simulation and manipulation aim to appropriate the knowledge of others by disinheriting them of their best judgments, creating a cause‑and‑effect relationship between the prior marauding and the orchestrated simulation.

It is a reading of symbols and signs that continues through synchronic and diachronic cycles. Intelligence is the tool that allows us to distinguish and differentiate the material world of reality, in which we live, from the immaterial world of imagination.

Without intelligence, only reality exists. Everything that belongs to knowledge and understanding—our ability to reason a priori about concepts, questions, situations, circumstances, etc.—does not exist.

For these individuals, the imaginary and all immaterial things are “the same as” material reality. Every fiction is assimilated to reality, and they accuse everything fictional of being real, without self‑criticism or remorse. A convenient way to perpetually harass and manipulate reality as they please.

\subsection{Depersonalization and Deviant Behaviors}

The individual suffering from idioschizophrenia seeks to be the person they are, but only in their thoughts. This person they attempt to become in their mind, they seek to feel through a second person who is not themselves.

From this depersonalization arise toxic, deviant, delinquent, or socially improper behaviors.

\subsection{The Fantasy of Enslavement}

The term “esclavagisme” is not accurate. The proper form is “esclavagiste,” which designates the person who engages in the trafficking of slaves. “Esclavage” is what the slave is subjected to.

There is no logical chain that leads someone to become a slave; this is why “esclavagisme” is not an expression in itself. “Esclavagisme” is more a fantasy than a reality.

The individual suffering from idioschizophrenia nourishes this fantasy intensely. They spend their time observing good conduct, particularly that of men. They watch for a term, a word, an expression, a turn of phrase that would allow them, through these elements, to subject these observed men to the slavery of a life of misery, misfortune, poverty, sadness, illness, death, etc.—forever.

This is the pernicious wish of the individual suffering from this condition: a secret handicap they wish to see imposed upon humanity in general.

This individual is misanthropic, meaning they harbor a general hatred of human beings. To illustrate this, the author imagines it as wanting to place one’s head inside one’s own mouth—a mystical image of hatred. This hatred is associated with the figure of death, sometimes represented with a large mouth.

Finally, the author describes idioschizophrenia as a condition resulting from an inability to know without reproach. The afflicted individual seeks to disproportion what is known so that this knowledge ceases to be recognized. This intellectual handicap can be heard when speaking with someone who suffers from it. They proclaim themselves the “police” of what exists and what does not exist, convinced they possess a superior sense of reality.

Quickly, they divert the conversation by asserting that what their interlocutor says “does not exist”—what the author calls the first response: “It‑does‑not‑exist‑to‑say.”

Then, they attribute everything the other person says to this initial assertion—the second response: “It‑is‑all‑because‑of‑to‑say.”

Such a discussion produces a morbid staging, a frenzy, in which the interlocutor feels discomfort. The afflicted person works to maraud the other’s opinion, to restrict the factors guiding their decisions and actions. They observe this process meticulously to provoke failures in the quality of the interlocutor’s judgment. The author calls this “frenzy.”

Thus, these individuals orchestrate a real‑time simulation to provoke the judgment error of others, depriving them of their best discernment in order to bring them down to their own handicap. In certain academic environments, such practices of simulation and manipulation aim to appropriate the knowledge of others by disinheriting them of their best judgments, creating a cause‑and‑effect relationship between the prior marauding and the orchestrated simulation.

\section{Consequences of Self‑Proclaimed Communication}

The subsidiarity that led Philippe Thomas Savard to act as guardian for individuals suffering from idioschizophrenia exposed him to a multitude of unfortunate events that brought him from the misery of homelessness to the brink of death twice. Twice, individuals suffering from idioschizophrenia came very close to causing Philippe’s death.

Both times, a succession of events unfolded in which these afflicted individuals orchestrated their usual frenzy, fabricating an environment in which they simulated within that same environment to provoke fault and, of course, to disinherit the author of his best judgment in order to steal or nullify knowledge. The marauding they carried out toward the people around Philippe—regarding work, friends, love, family, etc.—was omnipresent.

“The author’s sexuality has been hindered throughout the 41 years of his life during which he was placed in a position of guardianship toward these individuals. The level of sexual harassment, in addition to the assaults and brutality these individuals committed against Philippe, is daily and will end only with a single option they barely pronounce: ‘Negotiate.’ It is the only comment they make when the author complains about the situation, and the reason for this self‑referential evidence: ‘If you don’t know, it doesn’t matter, you don’t know,’ or ‘You’re not able to tell us what it does, so it does nothing.’”

The pseudo‑researcher of the curious study invokes politeness while the violation of Philippe’s privacy is omnipresent, alleging that the author must commit a sexual crime; otherwise, his conduct would demonstrate that not all men are rapists. She then asserts: “I do not agree.” The researcher always presents herself as a convinced homosexual who does not want women to benefit from themselves.

To demonstrate the antinomy of the reasoning and its self‑referential nature, the reason Philippe must commit this sexual crime, according to her, is that he must “assume it.” These reasons are equally unenviable: “so that this researcher may eternally have the trust of women in order to take advantage of them afterward.” This reasoning, which forces men into sexual delinquency, must apply to all men—except those who already have a history of sexual delinquency, according to her.

As a consequence of the obvious fact that these individuals are stubborn and significantly persuasive in depriving men in general of their exaggerated success in having a woman in their life.

This summary of the average worldview of the individual suffering from idioschizophrenia reveals much about the dystopian nature of their reasoning. Dystopian is, in the present, something that is certainly a utopia: this reasoning bases its assertions on something rather dubious: “Men with a history of sexual delinquency.”

To conclude, they attempt to enlighten us on the supposed “joy of life” we would have if we applied their reasoning: “As a consequence of the obvious fact that these individuals are stubborn and significantly persuasive in depriving men in general of their exaggerated success in having a woman in their life.” Dystopia, paranoid schizophrenia!

The word idioschizophrenia is a neologism formed by the author, Philippe Thomas Savard, from the words idiom and idiosyncratic.

Idiosyncratic: To succeed at something by applying oneself to something else. It is the success of the one who is conscious. When someone is conscious of what they are applying themselves to in the present moment, they achieve something unexpected, but they must be aware of what they are doing.

Idiom: An expression or term localized to a given region; outside that region, the idiom is not used or is used very little. Language allows the idiom, not the idiom that teaches language how to speak.

Idioschizophrenia: Early disturbance of the relationship between man and learned man. Depersonalization in which the individual seeks to be the person they are in their thoughts; this person is desired and felt through another person who is not themselves. A chirality from which deviant, delinquent, toxic, and improper behaviors arise. A delusion in which the rupture between the tangible world and imagination is no longer differentiated; self‑referential discourse, antinomy. Antinomism, free spirit: “Stubborn delusion to disproportion what is known in order to nullify the value of that knowledge.” An individual often excommunicated, living for fornication. A misanthropic individual who hates humanity in general. Placing one’s head in one’s mouth as the illustration of hatred and its disproportionate ambition toward knowledge. To illustrate the early nature of this disproportion: “Wanting to witness the premature abortion of one’s mother in order to swallow oneself.” Death is the equal of hatred and is often illustrated with a large mouth. A condition geographically localized to “Doctor.”

Doctus cum libro: Said of those incapable of thinking for themselves. They display borrowed knowledge and draw their ideas from the works of others. “Def. Pierre Larousse Dictionary, pink pages.” (Learned with the book).

Finesse, life drive: circumstantial fantasy that surpasses life through its reasons for being. The life drive, or finesse, is the particularity that every complex living being possesses within itself and which operates in the background in relation to the position it occupies in its environment. This position, which influences in the background or unconsciously, is that every individual possesses within themselves enough to always surpass a critical situation in their existence.

Consequence of autonomy, well‑being, and organization as the foundation of the structure of a living being. This drive, despite its vital importance, can lead to dark thoughts, suicide, and finality.

The subversion carried against autonomy can, like wartime pamphlets distributed to demotivate populations until they cease defending themselves, violate autonomy through this channel. As a consequence of the life drive operating in the background in a sick person, the best way to gain the upper hand over themselves would be to take their own life.

Their discourse often goes in the direction that it would be better for their loved ones if they left the living for the dead, and that even their loved ones understand this. It is a consequence of this finesse drive, which is more than indispensable but perceptible only by observing the organization of a living being.

The one who organizes themselves is meant to fully enjoy it. When there is external intervention, like a cell in autopoiesis, time stops so that the cell can regenerate within itself.

The vision of the person suffering an attack on the life drive, due to the distortion that subversion against autonomy causes, makes them see their loved ones as victims of their fate, whereas the reality is that they themselves are victims of a phenomenon that renders them incapable of realizing the gravity of their situation.

Friendship, which is the recognition that one does not have to do everything alone and that it is permissible to receive a helping hand from another, is overshadowed by opinions that seek to sever friendship by conveying that we must do everything alone. Opinions of fascist economies that advocate autarky. Alter‑globalization, the fascism of the new millennium, burdens the discussion to pollute it and, through subversion, put an end to the original.

\section*{Other Definitions}

\subsection*{Three Dominant Psychologies in a Person}

\textbf{Preschool Psychology}: A psychology observable at two different moments and in two different groups of people.

The first group adopting this preschool psychology is, of course, those who are logically associated with it by the sequence of events. These are the youngest members of our society who have not yet reached the level of first grade.

At this age, preschool psychology is one of acceptance and tolerance toward the various situations surrounding them. Indeed, at this age, they are neither tolerated nor considered for their decisions, still too young to make decisions for their future. Usually, their parents take charge of this and provide for their needs.

Social and moral customs grant parents the privilege of education and the responsibility of ensuring the needs of their offspring. Reality and common sense, in my view, suggest that every person has this responsibility toward children and adolescents when the occasion and circumstance require it. This is only an opinion, but I believe it may even be legally obligatory. We cannot ignore or fail to report physical or sexual violence against a minor in Canada.

This psychology is also observable when our elders lose autonomy and once again become more distant from the decision‑making that surrounds them. Circumstantial to the situation linked to their loss of autonomy. This psychology is again one of acceptance and self‑satisfaction toward the various situations of life.

\section{The Analogist Spirit and the Spirit of Finesse: A Narrated Reading}

\subsection{Introduction}

In the analogist perspective, understanding the world is not a matter of stacking rules, but of recognizing the hidden correspondences that link phenomena together. The spirit of finesse is the ability to perceive these connections even before demonstrating them. Together, these two faculties form a method: a way of seeing, sensing, and constructing knowledge.

\subsection{The Spirit of Finesse: An Inner Map of the Real}

The spirit of finesse is an elevation of oneself in relation to the present situation. Every human being evolves within an \emph{envelope}—a mental and concrete space where their objects, emotions, gestures, and habits reside. This ensemble forms an \textbf{inner map}, a living topology of our existence.

This map is not abstract: it is made of the position of our belongings, our actions, our emotions, our memories. It is the intimate geometry of our balance. The spirit of finesse allows us to read this map, to perceive what, in our immediate environment, can respond to a difficulty or illuminate a question.

\subsection{\textit{Lalangue}: When Words Touch the Body}

\textit{Lalangue} does not act through the definition of words, but through their \emph{physical impact}.  
Words leave traces, like pressures on the body. They reach the interlocutor from within.

To illustrate this, the analogist uses the \textbf{Klein bottle}: an impossible figure in which inside and outside merge.  
In the same way, \textit{lalangue} passes through the person: what is spoken outside folds back inside, without boundary.

This dynamic recalls \textbf{teleosemantics}, a branch of the philosophy of mind that seeks to map mental representations. The mind projects its forms outward, as if part of our neural network were reflected in the world.

\subsection{The Neural Network and Its Mirror in Society}

The human brain is a vast interconnected network.  
The analogist observes that certain zones of this network can project themselves into behavior, into the environment, into social interactions.

Analogously, the web is also a neural network:
\begin{itemize}
    \item it possesses memory,
    \item it stores information,
    \item it reacts to human behavior.
\end{itemize}

Thus, the analogist considers that certain behaviors—particularly \textbf{algorithmic biases}, meaning deceptive or harmful ways of acting—can affect both the human network and the digital network.

\subsection{The Role of the Analogist: Removing Algorithmic Biases}

An analogist is a grammarian who sees grammar as a form of mathematics.  
They assign value to words, gestures, and behaviors.  
When a behavior is recognized as deceptive, fraudulent, or harmful to the collective, it must be removed—just as one removes an incorrect term from an equation.

The analogist then acts like an electrical transformer:
\begin{itemize}
    \item the human neural network is the primary coil,
    \item the digital network is the secondary coil.
\end{itemize}

What circulates in one induces effects in the other.

\subsection{Recording Harmful Behaviors}

Certain individuals, often with limited schooling but close to academic environments, adopt behaviors analogous to algorithmic biases. They seek to disinherit the knowledge of others, to nullify the value of what is known, by repeating formulas such as:
\begin{itemize}
    \item “It‑does‑not‑exist‑to‑say,”
    \item “It‑is‑all‑because‑of‑to‑say.”
\end{itemize}

These individuals accept being recorded, believing they can escape consequences through fraud or manipulation. They even interpret explicit refusal from others as implicit consent, despite the legal fact that \textbf{no means no}.

\subsection{Machination and the Erroneous Belief in an Imaginary Dialogue}

When they are recorded, these individuals believe they are speaking to real people.  
They converse with an electronic device—sometimes equipped with a synthetic voice—but are convinced they are interacting with their victims.

They:
\begin{itemize}
    \item lurk around their targets,
    \item perform gestures without explanation,
    \item return to the device to “question” their own actions,
    \item justify their behaviors through imaginary connections.
\end{itemize}

This machination affects society concretely:
\begin{itemize}
    \item it disrupts the social environment,
    \item it introduces confidential information into the digital network,
    \item it unconsciously influences the behavior of everyone.
\end{itemize}

\subsection{The Feedback Effect: The Unconscious as an Alarm}

By acting in this way, their unconscious eventually informs others of their misconduct.  
Each targeted person intuitively senses the danger even before the act is committed.  
It is a natural mechanism of social self‑protection.

\subsection{Isossophy: The Analogist Method for Removing a Bias}

Isossophy is the method that allows the analogist to remove an algorithmic bias.  
It relies on an \textbf{equal measure} between:

\begin{itemize}
    \item real knowledge,
    \item and its deceptive excess.
\end{itemize}

It preserves the values of the past, protects those of the present, and prevents ignorance from being taught as truth.

\section*{Conclusion}

To all those who have taken the time to explore the work I have devoted to the geometry of the prime number spectrum, I wish to express my profound gratitude. Whether this project pleased you, puzzled you, or perhaps even disappointed you, I remain sincerely thankful that you offered a moment of your life to read a work I conceived entirely, from A to Z.

Thank you for your attention, your curiosity, and your encouragement. Thank you, above all, for having read me.

\bigskip

\noindent\textit{“There is no known scholar who does not remember the time when they were known for knowing. Those who no longer remember the time when they were recognized for their knowledge were never known, nor was their knowledge. Today, mine is—and I thank you for it.”}

\bigskip

To close this memoir, allow me a personal anecdote.  
Albert Einstein, when he presented relativity, was not a doctor. According to several biographies, he had a three‑year university education in industrial mechanics. Yet this training was enough for him to transform forever the school notebooks of the entire world.

For my part, I studied two years in civil engineering at the Centre intégré en mécanique industrielle de la Chaudière, in Saint‑Georges‑de‑Beauce. I then completed a one‑year training program in construction electricity.  
\[
2 + 1 = 3
\]

Thus, since the industrial mechanics training of the 1800s was sufficient for Einstein to conceive relativity and become one of the most influential scholars in history, I certainly do not claim to be as great a scholar as he was — but I, too, possess three years of study in this field. I therefore consider myself sufficiently equipped to address the subjects to which I allude in this work.

\bigskip

\begin{center}
\textbf{Thank you.}
\end{center}

%=========================================================
%   APPENDIX: APACHE LICENSE 2.0
%=========================================================

\appendix
\section*{Apache License 2.0}
\addcontentsline{toc}{section}{Apache License 2.0}

\begin{small}
\begin{verbatim}
                                 Apache License
                           Version 2.0, January 2004
                        http://www.apache.org/licenses/

TERMS AND CONDITIONS FOR USE, REPRODUCTION, AND DISTRIBUTION

1. Definitions.

"License" shall mean the terms and conditions for use, reproduction,
and distribution as defined by Sections 1 through 9 of this document.

"Licensor" shall mean the copyright owner or entity authorized by
the copyright owner that is granting the License.

"Legal Entity" shall mean the union of the acting entity and all
other entities that control, are controlled by, or are under common
control with that entity.

"Source" form shall mean the preferred form for making modifications,
including but not limited to software source code, documentation
source, and configuration files.

"Object" form shall mean any form resulting from mechanical
transformation or translation of a Source form.

"Work" shall mean the work of authorship, whether in Source or Object
form, made available under the License.

"Derivative Works" shall mean any work, whether in Source or Object
form, that is based on (or derived from) the Work and for which the
editorial revisions, annotations, elaborations, or other modifications
represent, as a whole, an original work of authorship.

"Contribution" shall mean any work of authorship intentionally
submitted to Licensor for inclusion in the Work by the copyright
owner or by an individual or Legal Entity authorized to submit on
behalf of the copyright owner.

"Contributor" shall mean Licensor and any individual or Legal Entity
on behalf of whom a Contribution has been received by Licensor.

2. Grant of Copyright License.
Subject to the terms and conditions of this License, each Contributor
hereby grants to You a perpetual, worldwide, non-exclusive, no-charge,
royalty-free, irrevocable copyright license to reproduce, prepare
Derivative Works of, publicly display, publicly perform, sublicense,
and distribute the Work and such Derivative Works in Source or Object
form.

3. Grant of Patent License.
(... full license text ...)
END OF TERMS AND CONDITIONS
\end{verbatim}
\end{small}

\end{document}